\chapter{El Secreto con Ruido: Aprendiendo con Errores (LWE)}
\label{cap:lwe_ruido}

\section{La aritmética del reloj: Números que dan vueltas}

Antes de comprender cómo funciona el problema más famoso de la criptografía moderna, necesitamos recordar una herramienta matemática que utilizas todos los días sin darte cuenta: \textbf{el reloj}.

Si tu reloj marca las $10$ de la mañana y tu clase dura $4$ horas, no dices que saldrás a las ``$14$ de la mañana'', sino a las $2$. ¿Qué hiciste mentalmente? Sumaste $10 + 4 = 14$, le restaste el ciclo completo de $12$, y te quedaste con el residuo:
\[
10 + 4 \equiv 2 \pmod{12}
\]

\begin{figure}[htbp]
\centering
\begin{tikzpicture}[scale=0.95]
  % Circulo
  \draw[line width=1.5pt, PQCMainBlue] (0,0) circle (2.2cm);
  
  % Horas 0 a 10 (Z_11)
  \foreach \i in {0,...,10} {
    \coordinate (H\i) at ({90 - \i*360/11}:2.2cm);
    \fill[PQCMainBlue] (H\i) circle (3pt);
    \node[font=\small\sffamily\bfseries] at ({90 - \i*360/11}:2.65cm) {\i};
  }
  
  % Flecha que ilustra dar la vuelta
  \draw[-{Stealth[length=3mm]}, line width=1.5pt, PQCTeal] (90:1.6cm) arc (90:-40:1.6cm);
  \node[font=\footnotesize\sffamily, align=center, PQCTeal] at (0.4, 0.5) {$8 + 5 = 13 \equiv 2$\\(en ciclo $q=11$)};
\end{tikzpicture}
\caption{La aritmética modular $\mathbb{Z}_{11}$: una rueda cerrada donde los números nunca escapan al infinito.}
\label{fig:reloj_modular}
\end{figure}

En criptografía, en lugar de trabajar con la recta real infinita $\mathbb{R}$ (donde los números con decimales crecen y consumen memoria infinita), fijamos un número entero positivo $q$ llamado \textbf{módulo}. 

El universo donde vivimos se llama $\mathbb{Z}_q = \{0, 1, 2, \dots, q-1\}$. Cada vez que una suma o multiplicación sobrepasa $q$, simplemente nos quedamos con el residuo de la división entre $q$.

\section{El problema LWE: El invento de Oded Regev}

En el año 2005, el matemático y científico de la computación Oded Regev formuló un problema computacional brillante llamado \textbf{LWE} (\textit{Learning With Errors}, o Aprendizaje con Errores). 

Su idea es tan elegante que cabe en una conversación de café:

\begin{intuicion}[¿Cómo se juega a LWE?]
\begin{enumerate}
  \item \textbf{El Secreto:} Imagina que tienes un vector secreto $\mathbf{s}$ compuesto por números enteros pequeños (por ejemplo $\mathbf{s} = (s_1, s_2)$).
  \item \textbf{Las Preguntas Públicas:} Alguien elige al azar una fila de números $\mathbf{a} = (a_1, a_2)$ a la vista de todo el mundo.
  \item \textbf{El Cálculo Limpio:} Haces el producto punto habitual:
  \[
  \mathbf{a} \cdot \mathbf{s} = a_1 s_1 + a_2 s_2 \pmod q
  \]
  \item \textbf{El Ruido Mágico:} En lugar de publicar ese resultado limpio, le sumas un error diminuto secreto $e$ (elegido al azar, por ejemplo de $\{-1, 0, 1\}$):
  \[
  b = \mathbf{a} \cdot \mathbf{s} + e \pmod q
  \]
\end{enumerate}
\end{intuicion}

Tú publicas los pares $(\mathbf{a}, b)$. Todo el mundo conoce las preguntas $\mathbf{a}$ y los resultados ruidosos $b$. 

\begin{ojo}[El desafío imposible para el atacante]
Dada una lista pública gigante de parejas $(\mathbf{a}_1, b_1), (\mathbf{a}_2, b_2), \dots, (\mathbf{a}_m, b_m)$, ¿puede un atacante descubrir el vector secreto $\mathbf{s}$?

\textbf{Respuesta:} ¡No! Si no hubiera error ($e=0$), el sistema se resolvería con el método de Gauss en milisegundos. Pero al existir el pequeño error desconocido $e$, el problema de hallar $\mathbf{s}$ es computacionalmente equivalente a resolver el problema del vector más cercano (CVP) en un retículo de alta dimensión.
\end{ojo}

\section{Un ejemplo numérico paso a paso}

Hagamos juntos un ejercicio real con números pequeños para ver la magia con tus propios ojos.

\begin{ejemplonumerico}[LWE con módulo $q=13$]
Elijamos un módulo pequeño $q = 13$.
Supongamos que nuestro secreto es un solo número: $s = 4$.

Generemos tres muestras públicas:
\begin{enumerate}
  \item \textbf{Muestra 1:} Elegimos al azar $a_1 = 3$. 
  Calculamos $3 \times 4 = 12$. Agregamos un pequeño ruido $e_1 = +1$:
  \[
  b_1 = 12 + 1 = 13 \equiv 0 \pmod{13}
  \]
  Publicamos el par $(3, 0)$.
  
  \item \textbf{Muestra 2:} Elegimos al azar $a_2 = 5$. 
  Calculamos $5 \times 4 = 20 \equiv 7 \pmod{13}$. Agregamos un ruido $e_2 = -1$:
  \[
  b_2 = 7 - 1 = 6 \pmod{13}
  \]
  Publicamos el par $(5, 6)$.
  
  \item \textbf{Muestra 3:} Elegimos al azar $a_3 = 2$. 
  Calculamos $2 \times 4 = 8$. Ruido $e_3 = 0$:
  \[
  b_3 = 8 + 0 = 8 \pmod{13}
  \]
  Publicamos el par $(2, 8)$.
\end{enumerate}

El público solo ve:
\[
(a_1=3, b_1=0), \quad (a_2=5, b_2=6), \quad (a_3=2, b_3=8)
\]
Para un atacante, los números $b$ parecen absolutamente aleatorios. No puede simplemente ``despejar'' $s$ porque no sabe qué porción de cada $b$ es información pura y qué porción es ruido.
\end{ejemplonumerico}

\section{¿Por qué una computadora cuántica no puede romper LWE?}

Las computadoras cuánticas son máquinas fascinantes, pero no son varitas mágicas capaces de probar todas las combinaciones al instante. 

El motivo por el cual una computadora cuántica puede resolver ciertos problemas tradicionales es que esos problemas antiguos poseen una \textbf{estructura periódica regular} (como una onda de radio limpia con una frecuencia exacta). El algoritmo cuántico de Shor es una especie de ``sintonizador de radio cuántico'' ultrasensible que detecta periodos regulares a la perfección.

\begin{intuicion}[La belleza del caos en LWE]
En el problema LWE, el ruido rompe cualquier patrón periódico. No hay una onda limpia que sintonizar: cada ecuación fue perturbada en una dirección diferente e impredecible. Por eso, ni las supercomputadoras clásicas más potentes ni las futuras computadoras cuánticas universales conocen un algoritmo eficiente para quebrar LWE en dimensiones altas.
\end{intuicion}

\begin{resumencapitulo}
\begin{itemize}
  \item La aritmética modular $\mathbb{Z}_q$ confina todos los cálculos a un anillo finito de números como las horas de un reloj.
  \item LWE formula ecuaciones lineales simples perturbadas por un pequeño error secreto: $\mathbf{b} = \mathbf{A}\mathbf{s} + \mathbf{e} \pmod q$.
  \item El ruido destruye la posibilidad de aplicar eliminación gaussiana, convirtiendo el problema en una búsqueda geométrica en un retículo.
  \item La ausencia de periodicidad regular protege a LWE tanto de atacantes convencionales como de computadoras cuánticas.
\end{itemize}
\end{resumencapitulo}
